% Detector to Diagnosis: LaTeX preamble additions (used with pandoc's default template).
% The design is in docs/PLAN-interior-design.md, section 3. The colors come from
% build/tex/palette-<profile>.tex (book/tools/palette.py), included before this file.
\usepackage[most]{tcolorbox}
\usetikzlibrary{calc,positioning,arrows.meta,shapes.geometric,shapes.misc,decorations.pathreplacing,fadings,shadows}
\usepackage{makeidx}
\makeindex
\usepackage{fvextra}
% Break long code lines instead of overflowing the margin.
\DefineVerbatimEnvironment{Highlighting}{Verbatim}{breaklines,breakanywhere,commandchars=\\\{\},fontsize=\small}
\fvset{breaklines=true,breakanywhere=true}
\usepackage{microtype}
\usepackage{needspace}
\usepackage{etoolbox}
\setlength{\emergencystretch}{3em}

% ---------------------------------------------------------------- the current Part
% \dsetpart{n} makes Part n's color the current one (partcolor): Part openers, chapter
% heads, running heads, section numbers, table heads, code rules and thumb tabs use it.
\newcommand{\dcurpart}{0}
\newcommand{\dtabpos}{0}
\newcommand{\dsetpart}[1]{\renewcommand{\dcurpart}{#1}\colorlet{partcolor}{d2dpart#1}%
  \ifstrequal{#1}{A}{\renewcommand{\dtabpos}{10}}{\ifstrequal{#1}{S}{\renewcommand{\dtabpos}{0}}{\renewcommand{\dtabpos}{#1}}}}
\newcommand{\dpartroman}[1]{\ifcase#1 \or I\or II\or III\or IV\or V\or VI\or VII\or VIII\or IX\fi}
% Bleed: on screen the Part band runs to the page edge; in the black ink interior (bw) it
% stays inside KDP's 0.125 in safety margin, since that interior has no bleed.
\newlength{\dbleed}
\ifdefstring{\dprofile}{bw}{\setlength{\dbleed}{0.3in}}{\setlength{\dbleed}{0pt}}

% ---------------------------------------------------------------- icons
% Small icons drawn in TikZ (no icon font), 1em high, in the current color.
\newcommand{\dicon}[1]{\begin{tikzpicture}[baseline=-0.32em,x=1em,y=1em,line width=0.08em,line cap=round,line join=round]\csname dicon@#1\endcsname\end{tikzpicture}}
\expandafter\def\csname dicon@why\endcsname{% a light bulb
  \draw (0,0.08) circle (0.32); \draw (-0.13,-0.22) -- (-0.13,-0.38) -- (0.13,-0.38) -- (0.13,-0.22);
  \draw (-0.1,-0.47) -- (0.1,-0.47); \draw (0,0.08) +(-0.12,0) .. controls +(0.06,0.14) .. +(0.12,0);}
\expandafter\def\csname dicon@regulated\endcsname{% a section sign in a shield
  \draw (-0.34,0.4) -- (0.34,0.4) -- (0.34,0.02) .. controls (0.34,-0.28) .. (0,-0.48) .. controls (-0.34,-0.28) .. (-0.34,0.02) -- cycle;
  \node[inner sep=0pt,font=\sffamily\bfseries\scriptsize] at (0,0.0) {\S};}
\expandafter\def\csname dicon@mistakes\endcsname{% a warning triangle
  \draw (0,0.44) -- (0.44,-0.4) -- (-0.44,-0.4) -- cycle; \draw[line width=0.1em] (0,0.18) -- (0,-0.12); \fill (0,-0.26) circle (0.05);}
\expandafter\def\csname dicon@exercise\endcsname{% a pencil
  \draw[rotate=45] (-0.12,-0.5) -- (-0.12,0.32) -- (0.12,0.32) -- (0.12,-0.5) -- (0,-0.7) -- cycle;
  \draw[rotate=45] (-0.12,0.18) -- (0.12,0.18);}
\expandafter\def\csname dicon@askteam\endcsname{% two speech bubbles
  \draw (-0.12,0.14) ellipse (0.34 and 0.24); \draw (-0.3,-0.06) -- (-0.4,-0.26) -- (-0.16,-0.1);
  \draw (0.2,-0.2) ellipse (0.26 and 0.18);}
\expandafter\def\csname dicon@milestone\endcsname{% a flag
  \draw (-0.3,-0.48) -- (-0.3,0.46); \draw (-0.3,0.44) -- (0.38,0.3) -- (0.12,0.15) -- (0.38,0.0) -- (-0.3,0.04);}
\expandafter\def\csname dicon@note\endcsname{% an information sign
  \draw (0,0) circle (0.42); \draw[line width=0.1em] (0,-0.22) -- (0,0.06); \fill (0,0.2) circle (0.05);}
\expandafter\def\csname dicon@check\endcsname{% a check mark in a box
  \draw (-0.4,-0.4) rectangle (0.4,0.4); \draw[line width=0.1em] (-0.22,0.0) -- (-0.05,-0.18) -- (0.24,0.2);}
\expandafter\def\csname dicon@summary\endcsname{% a list
  \foreach \y in {0.28,0,-0.28} {\fill (-0.34,\y) circle (0.05); \draw (-0.18,\y) -- (0.38,\y);}}
\expandafter\def\csname dicon@objectives\endcsname{% a target
  \draw (0,0) circle (0.42); \draw (0,0) circle (0.24); \fill (0,0) circle (0.07);}
\expandafter\def\csname dicon@prereq\endcsname{% a clock
  \draw (0,0) circle (0.42); \draw (0,0.26) -- (0,0) -- (0.18,-0.1);}

% ---------------------------------------------------------------- call-out boxes
% Every kind has a color, an icon and its own border, so it stays recognizable in gray.
% Titled boxes: a colored title bar with a white icon. #1 = semantic color, #2 = icon.
\tcbset{
  dbase/.style={enhanced,breakable,arc=1.5pt,boxrule=0pt,left=7pt,right=7pt,top=4pt,bottom=4pt,
    before skip=8pt,after skip=8pt,colback=#1!\dtint!white,colframe=#1,
    fonttitle=\sffamily\bfseries\small,coltitle=white,colbacktitle=#1,
    toptitle=1.5pt,bottomtitle=1.5pt,lefttitle=6pt,
    pad at break=2pt,parbox=false},
  dtitled/.style 2 args={dbase=#1,title={\dicon{#2}\hspace{0.45em}\dboxtitle}},
  dcard/.style 2 args={dbase=#1,leftrule=3pt,right=20pt,
    overlay unbroken and first={\node[anchor=north east,inner sep=3pt,text=#1] at (frame.north east) {\dicon{#2}};}},
}
\newcommand{\dboxtitle}{}
\newcommand{\dbox}[4]{% name, color, icon, extra options
  \newtcolorbox{#1}[1]{dtitled={#2}{#3},#4,code={\renewcommand{\dboxtitle}{##1}}}}
\dbox{boxwhy}{d2dwhy}{why}{leftrule=2.5pt}
\dbox{boxregulated}{d2dregulated}{regulated}{borderline west={1.1pt}{0pt}{d2dregulated},borderline west={1.1pt}{2.6pt}{d2dregulated},left=11pt}
\dbox{boxmistakes}{d2dmistakes}{mistakes}{borderline={0.9pt}{0pt}{d2dmistakes,dashed}}
\dbox{boxexercise}{d2dexercise}{exercise}{leftrule=2.5pt,colback=d2dexercise!\dtint!white}
\dbox{boxaskteam}{d2daskteam}{askteam}{borderline west={1.4pt}{0pt}{d2daskteam,dotted},left=10pt}
\dbox{boxmilestone}{d2dmilestone}{milestone}{boxrule=1.4pt}
\dbox{boxnote}{d2dnote}{note}{boxrule=0.4pt}
% Untitled cards for the recurring sections of a chapter.
\newtcolorbox{boxcheck}[1]{dcard={d2dcheck}{check}}
\newtcolorbox{boxsummary}[1]{dcard={d2dsummary}{summary}}
\newtcolorbox{boxobjectives}[1]{dcard={d2dobjectives}{objectives}}
\newtcolorbox{boxprereq}[1]{dcard={d2dnote}{prereq},colback=d2dnote!4!white}

% ---------------------------------------------------------------- code listings
% Pandoc defines Shaded from the highlight theme's background (templates/d2d.theme).
\renewenvironment{Shaded}{\begin{tcolorbox}[enhanced,breakable,colback=d2dcodebg,colframe=partcolor,
  boxrule=0pt,leftrule=2pt,arc=0pt,outer arc=0pt,left=5pt,right=4pt,top=2pt,bottom=2pt,
  before skip=4pt,after skip=8pt,pad at break=1pt]}{\end{tcolorbox}}
% The caption of an included listing (include-code.lua), kept with its listing.
\newcommand{\dlistingcap}[1]{\par\needspace{4\baselineskip}\noindent{\sffamily\footnotesize
  \textcolor{partcolor}{\rule[-0.15em]{2pt}{0.9em}}\hspace{0.4em}\textcolor{d2dmuted}{#1}}\par\nopagebreak\vspace{-2pt}}

% ---------------------------------------------------------------- tables
% A tinted head row (book/build.sh adds \rowcolor{dtablehead} to each longtable head) and
% light rows in alternation.
\colorlet{dzebra}{d2drule!22!white}
\AtBeginEnvironment{longtable}{\colorlet{dtablehead}{partcolor!\dstrongtint!white}\rowcolors{2}{white}{dzebra}\sffamily\small}
\setlength{\arrayrulewidth}{0.4pt}

% ---------------------------------------------------------------- figures
\usepackage{caption}
\captionsetup{font=small,labelfont={bf,sf,color=partcolor},textfont={color=d2dink}}
\captionsetup[figure]{name=Figure}

% ---------------------------------------------------------------- headings
\usepackage{titlesec}
\titleformat{\section}{\needspace{5\baselineskip}\Large\sffamily\bfseries\color{d2dink}}{\textcolor{partcolor}{\thesection}}{0.6em}{}
\titleformat{name=\section,numberless}{\needspace{5\baselineskip}\Large\sffamily\bfseries\color{d2dink}}{}{0pt}{\textcolor{partcolor}{\rule[-0.1em]{3pt}{0.95em}}\hspace{0.5em}}
\titleformat{\subsection}{\large\sffamily\bfseries\color{d2dink}}{\textcolor{partcolor}{\thesubsection}}{0.6em}{}
\titleformat{\subsubsection}{\normalsize\sffamily\bfseries\color{d2dink}}{\thesubsubsection}{0.6em}{}
\titlespacing*{\section}{0pt}{16pt plus 4pt minus 2pt}{6pt plus 2pt}
\titlespacing*{\subsection}{0pt}{12pt plus 3pt minus 2pt}{4pt plus 2pt}

% The chapter head: a tinted band with the number in a Part-colored badge, the title, and the
% chapter's emblem (book/figures/emblems.tex) on the right.
\newcommand{\demblemid}{}
\newcommand{\dsetemblem}[1]{\renewcommand{\demblemid}{#1}}
\input{emblems.tex}
\newcommand{\dchapteremblem}{%
  \ifcsname demblem@\demblemid\endcsname\csname demblem@\demblemid\endcsname
  \else\csname demblem@default\endcsname\fi}
\newcommand{\dchapterhead}[2]{% #1 = label (number or empty), #2 = title
  \noindent\begin{tikzpicture}
    \def\W{\textwidth}\def\H{1.32in}
    \fill[partcolor!\dtint!white,rounded corners=3pt] (0,0) rectangle (\W,\H);
    \fill[partcolor,rounded corners=3pt] (0,0) rectangle (1.05in,\H);
    \ifx\relax#1\relax
      \node[text=white] at (0.525in,0.5*\H) {\scalebox{2.2}{\dicon{summary}}};
    \else
      \node[text=white,font=\sffamily\bfseries\scriptsize,anchor=north] at (0.525in,\H-0.12in) {\addfontfeature{LetterSpace=12}CHAPTER};
      \node[text=white,font=\sffamily\bfseries\fontsize{40}{40}\selectfont,anchor=center] at (0.525in,0.43*\H) {#1};
    \fi
    \node[anchor=west,text width=\W-1.05in-1.55in,font=\sffamily\bfseries\LARGE,text=d2dink,align=left]
      at (1.25in,0.5*\H) {\raggedright #2};
    \begin{scope}[shift={(\W-0.72in,0.5*\H)}]
      \begin{scope}[x=0.58in,y=0.58in]\dchapteremblem\end{scope}
    \end{scope}
  \end{tikzpicture}\par}
\titleformat{\chapter}[block]{}{}{0pt}{\dchapterheadnum}
\titleformat{name=\chapter,numberless}[block]{}{}{0pt}{\dchapterhead{}}
\newcommand{\dchapterheadnum}[1]{\dchapterhead{\thechapter}{#1}}
\titlespacing*{\chapter}{0pt}{-24pt}{18pt}

% ---------------------------------------------------------------- Part openers
% \dpart{n}{roman}{title}{lit stages}{chapters}: a full page in the Part color with the
% MiniDAQ map (book/figures/tikz/pipeline-map.tex), the stages this Part works on lit, and
% the Part's chapters (\dpartch{number}{title}{label}), linked, with their pages.
\newcommand{\dpartch}[3]{\noindent\hyperref[#3]{\colorbox{partcolor}{\makebox[1.7em]{\color{white}\sffamily\bfseries #1}}%
  \hspace{0.8em}{\sffamily\bfseries\color{d2dink}#2}}\dotfill{\sffamily\color{d2dmuted}\pageref*{#3}}\par\vspace{5pt}}
\input{pipeline-map.tex}
\newcommand{\dpart}[5]{%
  \clearpage\dsetpart{#1}\thispagestyle{empty}%
  \phantomsection\addcontentsline{toc}{part}{Part #2: #3}%
  \markboth{}{}%
  \dprepmap{#4}{\dimexpr\paperwidth-1.7in\relax}%
  \begin{tikzpicture}[remember picture,overlay]
    \coordinate (nw) at ($(current page.north west)+(\dbleed,-\dbleed)$);
    \coordinate (ne) at ($(current page.north east)+(-\dbleed,-\dbleed)$);
    \coordinate (band) at ($(current page.north west)!0.5!(current page.south west)$);
    \shade[top color=partcolor!82!black,bottom color=partcolor] (nw) rectangle ($(ne |- band)$);
    % a ring of detector blocks, as on the cover
    \begin{scope}[shift={($(ne |- band)+(-1.6in,2.4in)$)},white,opacity=0.16]
      \foreach \a in {0,10,...,350} {\fill[rotate=\a] (2.2in,-0.09in) rectangle (2.5in,0.09in);}
      \foreach \a in {5,15,...,355} {\fill[rotate=\a] (2.55in,-0.07in) rectangle (2.75in,0.07in);}
    \end{scope}
    \node[anchor=north west,text=white,font=\sffamily\bfseries\small] at ($(current page.north west)+(1in,-1.15in)$) {\addfontfeature{LetterSpace=25}PART};
    \node[anchor=north west,text=white,font=\sffamily\bfseries\fontsize{96}{96}\selectfont] at ($(current page.north west)+(0.94in,-1.4in)$) {#2};
    \node[anchor=south west,text width=\paperwidth-2in,text=white,font=\sffamily\bfseries\fontsize{26}{31}\selectfont,align=left]
      at ($(band)+(1in,0.45in)$) {\raggedright #3};
    \node[anchor=north,inner sep=0pt] at ($(band)+(0.5\paperwidth,-0.6in)$) {\usebox{\dmapbox}};
    \node[anchor=north,font=\sffamily\small,text=d2dmuted] at ($(band)+(0.5\paperwidth,-0.6in-\dmapheight-0.12in)$) {Where this Part works in MiniDAQ: the stages in color.};
    \node[anchor=north west,text width=\paperwidth-2.4in,font=\normalsize]
      at ($(band)+(1.2in,-0.6in-\dmapheight-0.6in)$) {#5};
  \end{tikzpicture}%
  \clearpage}

% ---------------------------------------------------------------- running heads and tabs
\usepackage{fancyhdr}
\pagestyle{fancy}
\fancyhf{}
\fancyhead[LE]{\sffamily\footnotesize\textcolor{d2dmuted}{\addfontfeature{LetterSpace=8}DETECTOR TO DIAGNOSIS}}
\fancyhead[RO]{\sffamily\footnotesize\textcolor{d2dmuted}{\nouppercase{\leftmark}}}
\fancyfoot[LE,RO]{\sffamily\small\bfseries\textcolor{partcolor}{\thepage}}
\renewcommand{\headrule}{{\color{partcolor}\hrule height 0.6pt\vspace{-0.6pt}}}
\fancypagestyle{plain}{\fancyhf{}\fancyfoot[LE,RO]{\sffamily\small\bfseries\textcolor{partcolor}{\thepage}}\renewcommand{\headrulewidth}{0pt}\renewcommand{\headrule}{}}
% A thumb tab on the outer edge of every page of a Part, one step lower per Part, so a Part
% can be found by flipping. Only on screen and on the tablet: the KDP interior has no bleed.
\newcommand{\dthumbtab}{%
  \ifnum\dtabpos>0 \ifdefstring{\dprofile}{bw}{}{%
    \begin{tikzpicture}[remember picture,overlay]
      \pgfmathsetmacro{\dtabstep}{(\paperheight/72.27-1.8)/10}
      \pgfmathsetmacro{\dtaby}{0.9+(\dtabpos-1)*\dtabstep}
      \pgfmathsetmacro{\dtabh}{0.88*\dtabstep}
      \ifodd\value{page}
        \fill[partcolor] ($(current page.north east)+(-0.14in,-\dtaby in)$) rectangle ++(0.14in,-\dtabh in);
      \else
        \fill[partcolor] ($(current page.north west)+(0,-\dtaby in)$) rectangle ++(0.14in,-\dtabh in);
      \fi
    \end{tikzpicture}}\fi}
\makeatletter\AddToHook{shipout/background}{\if@mainmatter\put(0,0){\dthumbtab}\fi}\makeatother
